\subsection{Strict curvature and local relative robustness}
\label{app:stricttransitionrobustness}
The complete-graph example need not rely on zero curvature in terminal
energy or exact saturation of the added edges' post-early capacity.
The following strengthening remains an analytical compulsory-task model;
it is not an implementation qualification or a perturbation certificate
for the public timetable.

\begin{proposition}[A strictly convex, locally robust obstruction]
\label{prop:stricttransitionrobustness}
In Proposition~\ref{prop:completegraphtransition}, keep the within-block
edges and put $M=S_{\max}$, $\eta=1/4$. Replace each interblock task
energy by $d_{ij}=B-b_j-\eta$, and increase the terminal power/budget to
$H=2kB+1$. Choose per-bus acceptance power strictly larger than
$\max\{K,H\}$ at the nominal point, for example $\max\{K,H\}+1$,
and hold that acceptance power fixed locally. Here $\eta$ is an energy
margin, not a conversion efficiency.
Every perfect matching remains feasible. The coordinatewise strictly
increasing, strongly convex supply objective
\[
 F(E,L)=3E+L+\frac{(E-3)^2}{2K}
             +\frac{(L-(M-9))^2}{4M}
\]
gives the whole-problem values
\[
 \CH=M+2kf,\qquad
 D-\CH=\frac{9}{2K}+\frac{81}{4M}>0.
\]
For each fixed finite family member, positive gap and feasibility of every
perfect matching persist in a sufficiently small relative open
neighborhood of its numerical parameters, with topology, timing, access,
fleet cap and charge-only/unit-efficiency law fixed. Economic perturbations
are restricted to positive separable quadratics with positive coordinate
derivatives at the origin. No neighborhood radius or unchanged optimizer
is asserted.
\end{proposition}

\begin{proof}
For an added edge, post-early consumption is $B-\eta<B$ and its early
requirement is $r_{ij}=a_i-\eta>0$. Its early interval
$[r_{ij},a_i]$ has width $\eta$. All within-block intervals also have
strictly positive widths and strictly feasible post-early consumption
under $B\ge20w_{\max}$. Thus every matching has $R<K$; moreover
$S-R=\sum\min(e_{ij},B)\le2kB<H$. Serial common-window charging
realizes its entire branch $R\le E\le K$, $L=S-E$. Interior allocations
can be chosen strictly inside every per-bus early interval, leaving
early-budget and preterminal SOC margins as well as terminal-budget
slack. Final inventory remains exactly full by the subsequent refill.

The potentials in Proposition~\ref{prop:completegraphtransition} remain
tight on within-block edges. On an added edge, $e_{ij}+2r_{ij}$ drops by
exactly $3\eta$; its potential slack therefore remains at least
$3-3/4=9/4$. Every physical point and every complete-plan mixture obey
$3E+L\ge M$, while a physical plan using an interblock edge obeys
$3E+L\ge M+9/4$.

The Hessian of $F$ is $\operatorname{diag}(1/K,1/(2M))$. Its coordinate
derivatives on the nonnegative orthant are bounded below by
$3-3/K>0$ and $1/2+9/(2M)>0$. The equal mixture of the retained
corners $R=0$ and $R=6$ gives $(3,M-9)$ and makes both squares zero. Together with
the affine bound this proves $\CH=M+2kf$, uniquely in projected loads.
On a within-block branch $L=M-2R-E$, differentiation gives
\[
 \frac{dF}{dE}
 =2+\frac{E-3}{K}-\frac{L-(M-9)}{2M}
 \ge2-\frac3K-\frac9{2M}>0.
\]
Its minimum is at $E=R$, with excess
$(R-3)^2(1/(2K)+9/(4M))$. Exactly $R=0,6$ minimize it. Their excess is
at most $9/28+81/160<1<9/4$, using $K\ge14$, $M\ge40$.
No interblock plan can improve or tie them, so the stated physical
value covers every matching of the complete graph.

For the neighborhood statement, fix this finite labeled model. Allow
continuous numerical changes within the admissible nonnegative-energy,
nonnegative-fee parameter space, keeping the declared access and
serialization conditions. The strict acceptance margin persists for
sufficiently small changes of $K,H$, so fixed acceptance still permits
serialization. Zero task energies and $f=0$ are interpreted
in the relative topology. If identities such as $K=\sum_i a_i$ are
retained, use the relative topology of that subspace. Strict per-edge
interval widths, post-early capacity margins and terminal slack persist
simultaneously because only finitely many edges and matchings occur.

For independently varying budgets a matching's branch endpoints are
\[
 \ell_M=\max\{0,R_M,S_M-H\},\qquad
 u_M=\min\{\textstyle\sum_i a_i,K\}.
\]
These are continuous, including when
$R_M=\sum(e_{ij}-B)_+$ changes its active positive parts, and remain
strictly ordered locally. Parameterize each interval by $t\in[0,1]$.
The resulting continuous objectives on a common compact domain have
continuous minima; taking the minimum over finitely many matchings
proves continuity of $D$. Their complete-plan hull is generated by the
finitely many continuous branch endpoints. Parameterizing it by a fixed
simplex proves continuity of $\CH$ in the same way, with no
nondegeneracy assumption on hull faces. All nearby loads lie in a common
bounded box. The nominal strict gap therefore stays positive locally.
Positive separable curvature and positive coordinate derivatives at the
origin also persist locally, preserving strong convexity and global
monotonicity on the nonnegative orthant.
\end{proof}

This continuity argument does not preserve the exact displayed gap
formula, its two minimizing corners or its particular hull mixture after
perturbation. It gives no radius uniform in $k$, no robustness to new
physics or access options, and no assertion for arbitrary quadratic
cross-coefficients. In particular, it cannot transfer the public
binary-energy certificate to a different consumption law: the explicit
strict feasibility margins and finite-branch description used here must
first be proved for the model in question. Neither practical prevalence,
road realization nor computational hardness follows.
